\SourceLeaf{356}%
% A544-P0356-D01
\Article{5}
Hieraus leiten wir folgende zweite Methode den Decidenten zu bestimmen ab.

I. Man sammle alle Combinationen von ganzen Werthen von $Y$ und $x$, die folgende Eigenschaften haben
\begin{enumerate}
\item dass $\xi=\frac{Ax+bY}{\alpha}$ zwischen $0$ und $\frac12$ falle, wobei die zweite Grenze inclusive genommen wird
\item dass $\eta=\frac{-Bx+aY}{\alpha}$ zwischen $0$ und $\frac12$ falle, die erste Grenze inclusive genommen.
\end{enumerate}

% A544-P0356-D02
II. Man berechne dafür
\[
\begin{aligned}
X&=\frac{Dx+\beta Y}{\alpha}\\[3pt]
y&=\frac{\beta x+dY}{\alpha}
\end{aligned}
\]

% A544-P0356-D03
III. Man lasse alle diejenigen weg, wo $[X]$ eine ungerade Zahl ist, und theile die übrigen, wo $[X]$ gerade ist, in zwei Classen;

in die erste Classe setze man diejenigen, wo zugleich
\[
\begin{array}{l}
Y\text{ gerade, }x\text{ gerade, }[y]\text{ gerade}\\
\text{oder wo eine dieser Bedingungen Statt findet;}
\end{array}
\]
in die zweite Classe diejenigen, wo zwei dieser Bedingungen oder gar keine Statt hat,
\[
\begin{array}{rl}
\text{oder in I.}&\text{wo }[Y+x+y]\text{ gerade}\\
\text{II.}&\text{wo }[Y+x+y]\text{ ungerade}
\end{array}
\]
und nenne den Ueberschuss der Anzahl in der ersten Classe über die in der zweiten $c$.

% A544-P0356-D04
IV. Man sammle alle Combinationen von ganzen Werthen von $Y$ und $y$, die folgende Eigenschaften haben:
\begin{enumerate}
\item dass $\xi=\frac{Ay-aY}{\beta}$ zwischen $0$ und $\frac12$ falle, die erste Grenze inclusive, wenn $\beta$ negativ, die zweite inclusive, wenn $\beta$ positiv;
\item dass $\eta=\frac{-By+bY}{\beta}$ zwischen $0$ und $\frac12$ falle, die erste Grenze inclusive bei positivem $\beta$, die zweite bei negativem.
\end{enumerate}
\BookRule{.3\textwidth}

\SourceLeaf{357}%
% A544-P0357-D01
\par\bigskip\Needspace{7\baselineskip}
\pdfbookmark[1]{Zur Theorie der biquadratischen Reste. V.}{nachlass-biquadratic-V}
\centerline{[V.]}\medskip
\renewcommand{\Article}[1]{\par\medskip\Needspace{4\baselineskip}\pdfbookmark[2]{Art. #1}{nachlass-biquadratic-V-art-#1}\centerline{\textbf{[#1.]}}\nobreak\smallskip}
\Article{1}
\[
\begin{array}{lll}
\text{Modulus}&M=A+Bi&AA+BB=D\\
\text{Rest}&m=a+bi,&aa+bb=d
\end{array}
\]
\[
\frac{mD}{M}=\mu=\alpha+\beta i=aA+bB+(Ab-Ba)i
\]
\[
\xi+\eta i=\pi,\quad\pi m=p=x+yi;\quad\pi M=P=X+Yi
\]
$\omega$ eine unbestimmte unendlich kleine reelle positive Grösse.

% A544-P0357-D02
\Article{2}
\begin{center}\emph{Vorbereitung.}\end{center}
I. Man sammle alle $\pi$, wo
\[
\begin{array}{l}
\xi\text{ nicht negativ und nicht grösser als }\frac12\\
\eta\text{ positiv und kleiner als }\frac12\\
\text{Entweder }x\text{ oder }y\text{ eine Ganze}\\
\text{Entweder }X\text{ oder }Y\text{ eine Ganze}
\end{array}
\]
und bestimme für jedes $\pi$ die Grösse $\varepsilon$ nach folgender Regel:
\[
\begin{array}{ll}
\text{Es sei }p^0&\text{die nächste Ganze durch }1+i\text{ theilbare bei }p\\
P^0&\text{die nächste Ganze durch }1+i\text{ theilbare bei }P
\end{array}
\]
und setze $\varepsilon=\pm1$, wo das Zeichen immer dasselbe ist wie das Zeichen des imaginären Theils der Grösse
\[
\frac{p-p^0}{P-P^0}(\alpha-\beta i)
\]

% A544-P0357-D03
folgendes sind die Specialregeln: Erste Classe, $x$ und $X$ Ganze
\[
\begin{aligned}
\beta\xi&=-Bx+bX\\
\beta\eta&=-Ax+aX\\
\beta y&=-\alpha x+dX\\
\beta Y&=-Dx+\alpha X
\end{aligned}
\]

\SourceLeaf{358}%
% A544-P0358-D01
$\varepsilon=-1$, wenn $\beta$ positiv, $x$ gerade, $[y]$ gerade, $X$ gerade, $[Y]$ gerade oder wenn nur eine ungerade Anzahl dieser Bedingungen gilt.

$\varepsilon=+1$, wenn keine oder eine gerade Anzahl dieser Bedingungen gilt.

% A544-P0358-D02
Zweite Classe, $y$ und $Y$ Ganze
\[
\begin{aligned}
\beta\xi&=+Ay-aY\\
\beta\eta&=-By+bY\\
\beta x&=+\alpha y-dY\\
\beta X&=+Dy-\alpha Y
\end{aligned}
\]
$\varepsilon=-1$, wenn $\beta$ positiv, $[x]$ gerade, $y$ gerade, $[X]$ gerade, $Y$ gerade oder wenn eine ungerade Anzahl dieser Bedingungen gilt.

$\varepsilon=+1$, wenn keine oder eine gerade Anzahl gilt.

% A544-P0358-D03
Dritte Classe, $y$ und $X$ Ganze
\[
\begin{aligned}
\alpha\xi&=+By+aX\\
\alpha\eta&=+Ay-bX\\
\alpha x&=-\beta y+dX\\
\alpha Y&=+Dy-\beta X
\end{aligned}
\]
$\varepsilon=-1$, wenn $\alpha$ positiv, $[x],y,X,[Y]$ alle Gerade oder wenn eine ungerade Anzahl dieser Bedingungen gilt.

$\varepsilon=+1$, wenn keine oder eine gerade Anzahl Statt hat.

% A544-P0358-D04
Vierte Classe, $x$ und $Y$ Ganze
\[
\begin{aligned}
\alpha\xi&=+Ax+bY\\
\alpha\eta&=-Bx+aY\\
\alpha y&=+\beta x+dY\\
\alpha X&=+Dx+\beta Y
\end{aligned}
\]
$\varepsilon=+1$, wenn $\alpha$ positiv, $x,[y],[X],Y$ alle Gerade oder bei einer ungeraden Anzahl dieser Bedingungen.

$\varepsilon=-1$, bei keiner oder einer geraden Anzahl derselben.

\SourceLeaf{359}%
% A544-P0359-D01
\Article{3}
II. Man sammle alle $\pi$, wo
\[
\begin{array}{l}
\xi\text{ positiv und kleiner als }\frac12\\
\eta=\omega\\
\text{und entweder }X\text{ oder }Y\text{ eine Ganze,}
\end{array}
\]
und setze $\varepsilon=\pm1$ so dass das Zeichen des imaginären Theils von
\[
\frac{M}{P-P^0}
\]
zu nehmen ist.

Specialregel: Erste Classe, $X$ ganz
\[
\begin{aligned}
A\xi&=+X+B\omega\\
Ax&=+aX-\beta\omega\\
Ay&=+bX-\alpha\omega\\
AY&=+BX+D\omega
\end{aligned}
\]
$\varepsilon=-1$, wenn $A$ positiv, $X$ und $[Y]$ gerade oder wenn nur eine Bedingung gilt.

$\varepsilon=+1$, wenn keine oder zwei gelten.

% A544-P0359-D02
Zweite Classe, $Y$ ganz
\[
\begin{aligned}
B\xi&=+Y-A\omega\\
Bx&=+aY-\alpha\omega\\
By&=+bY-\beta\omega\\
BX&=+AY-D\omega
\end{aligned}
\]
$\varepsilon=+1$, wenn $B$ positiv, $[X]$ und $Y$ gerade oder wenn nur eine Bedingung gilt.

$\varepsilon=-1$, wenn keine oder zwei gelten.

% A544-P0359-D03
\Article{4}
III. Man sammle alle $\pi$, wo
\begin{quote}
$\xi$ und $\eta$ denselben Bedingungen unterworfen sind wie in II.\par
entweder $x$ oder $y$ Ganze,
\end{quote}
und setze $\varepsilon=\pm1$ mit dem Zeichen des imaginären Theils von
\[
\frac{m}{p-p^0}
\]

\SourceLeaf{360}%
% A544-P0360-D01
Specialregeln: Erste Classe, $x$ ganz
\[
\begin{aligned}
a\xi&=+x+b\omega\\
ay&=+bx+d\omega\\
aX&=+Ax+\beta\omega\\
aY&=+Bx+\alpha\omega
\end{aligned}
\]
$\varepsilon=-1$, wenn $a$ positiv, $x,[y]$ beide gerade oder wenn nur eine Bedingung gilt.

$\varepsilon=+1$, wenn keine oder zwei gelten.

% A544-P0360-D02
Zweite Classe, $y$ ganz
\[
\begin{aligned}
b\xi&=+y-a\omega\\
bx&=+ay-d\omega\\
bX&=+Ay-\alpha\omega\\
bY&=+By+\beta\omega
\end{aligned}
\]
$\varepsilon=+1$, wenn $b$ positiv, $[x]$ und $y$ gerade oder wenn nur eine Bedingung gilt.

$\varepsilon=-1$, wenn keine gilt.

% A544-P0360-D03
\Article{5}
IV. Man sammle alle $\pi$, wo
\[
\begin{array}{l}
\xi=\frac12+\frac12\omega\\
\eta\text{ positiv und kleiner als }\frac12\\
X\text{ oder }Y\text{ ganz,}
\end{array}
\]
und setze $\varepsilon=\pm1$ mit dem Zeichen des imaginären Theils von
\[
\frac{iM}{P-P^0}
\]
Specialregeln: Erste Classe $X$ eine Ganze,
\[
\begin{aligned}
2B\eta&=+A-2X+A\omega\\
2Bx&=-\beta+2bX-\beta\omega\\
2By&=+\alpha-2aX+\alpha\omega\\
2BY&=+D-2AX+D\omega
\end{aligned}
\]
$\varepsilon=+1$, wenn $B$ positiv, $X,[Y]$ gerade oder bei einer Bedingung,

$\varepsilon=-1$, bei keiner oder zwei Bedingungen.

